\cleardoublepage
\vspace*{\fill}
\begin{abstract}
    This thesis explores the engineering lifecycle of developing a Small Language Model (SLM) from scratch, utilizing a hybrid architecture of Transformers and Gated Recurrent Units (GRU). In a landscape dominated by massive models, the primary barrier for custom development is the extreme economic cost. This research shifts the focus from model scale to the complete deployment pipeline. We document the implementation of CUDA-optimized code, orchestration on AWS EC2 environments, and data strategies for conversational agents. Furthermore, we integrate Retrieval-Augmented Generation (RAG) to provide the model with external, up-to-date information, significantly reducing hallucinations without the prohibitive cost of full-scale retraining. Through performance and cost-benefit analysis using spot instances, we evaluate the feasibility of local development versus cloud scaling. Finally, we demonstrate production readiness by deploying the system on Kubernetes, visualized and managed through ArgoCD.
\end{abstract}
\vspace*{\fill}
\clearpage

\vspace*{\fill}
\begin{abstract}
    Este trabajo de fin de máster investiga el ciclo de vida de ingeniería en el desarrollo de un Small Language Model (SLM) desde cero, integrando arquitecturas Transformer y GRU. En un entorno donde el desarrollo de modelos de lenguaje está limitado por altos costes operativos, este proyecto se centra en el flujo completo de ingeniería más allá de la escala del modelo. La investigación abarca el desarrollo de código optimizado para CUDA, la gestión de infraestructura en AWS EC2 y la preparación de datos para agentes conversacionales. Se incluye además la implementación de tecnologías RAG (Retrieval-Augmented Generation), que permiten al modelo consultar fuentes externas de información para responder con precisión sin necesidad de reentrenamientos costosos. Mediante el uso de instancias spot y subconjuntos de datos, se analizan los costes y el rendimiento del sistema. Finalmente, se valida la puesta en producción mediante un despliegue en Kubernetes gestionado con ArgoCD.
\end{abstract}
\vspace*{\fill}
\clearpage

\vspace*{\fill}
\begin{abstract}
    Aquest treball de fi de màster investiga el cicle de vida d'enginyeria en el desenvolupament d'un Small Language Model (SLM) des de zero, integrant arquitectures Transformer i GRU. En un entorn on el desenvolupament de models de llenguatge està limitat per alts costos operatius, aquest projecte se centra en el flux complet d'enginyeria més enllà de l'escala del model. La investigación inclou el desenvolupament de codi optimitzat per a CUDA, la gestió d'infraestructura a AWS EC2 i la preparació de dades per a agents conversacionals. S'hi inclou a més la implementació de tecnologies RAG (Retrieval-Augmented Generation), que permeten al model consultar fonts externes d'informació per respondre amb precisió sense la necessitat de reentrenaments costosos. Mitjançant l'ús d'instàncies spot i subconjunts de dades, s'analitzen els costos i el rendiment del sistema. Finalment, es valida la posada en producció mitjançant un despliegament a Kubernetes gestionat amb ArgoCD.
\end{abstract}
\vspace*{\fill}
\clearpage
